\section{Conclusion}
\label{sec:conclusion}

We asked whether the $\sim$80\% of \sae features that do not align with any \swissprot concept in a public \plm encode previously unknown biology. On the pretrained \interplm \sae for \esm layer 4, evaluated against 422{,}146 residues of human \swissprot, the answer is largely \emph{no}: the top of the dark tail encodes recognizable sub-domain motifs (\tgekp linkers, \dry and \npxxy GPCR motifs) that the annotation system describes at coarser granularity, and adding a single concept (``Zinc finger'') re-claims 7 of 8 top features at F1 $\geq 0.6$. \sae features in this model regime are best understood as an annotation-refinement tool rather than a discovery engine. In parallel, our forward-hook ablation experiment provides a quantitative negative result for the ``F1 implies mechanism'' inferential step: none of the 30 top dark features pass single-feature ablation under Benjamini--Hochberg FDR, and the median KL effect is in fact $2^{-3.86}\times$ \emph{smaller} at activating residues than at matched non-activating controls.

\para{Key takeaway.} High feature--concept correlation, even with high LLM-rated biological confidence, does not imply mechanistic role under single-feature ablation in \esm. Dark features in this small \plm encode known sub-domain motifs that \swissprot describes at coarser granularity. Both claims hold at the scale of $10^4$ \sae features and $10^5$ residues we evaluated, with code that runs end-to-end in 25 minutes.

\para{Future work.} The natural follow-ups are (i) scaling the pipeline to \esmlarge and \texttt{InterPLM-esm2-650m} layers 9, 18, and 24, where features are less superposed and genuine novelty becomes more plausible; (ii) re-running ablation with Ordered/top-$k$ \sae checkpoints to test whether the manifold-robustness explanation is architecture-specific; (iii) co-ablating clusters of decoder-collinear features rather than single features, to test the distributed-encoding hypothesis; (iv) running the full 217-assay \proteingym pipeline to recover statistical power for the functional-correlation axis; and (v) annotating the full 1{,}207 structured dark features with \gptfive paired with literature retrieval, which is where genuinely post-\swissprot biology would surface if it exists. We release the pipeline and outputs to make these follow-ups concrete.
